\documentclass[12pt, a4paper]{exam}
\usepackage[a4paper,top=2.3cm,bottom=2.8cm,left=2cm,right=2cm,headheight=15pt,headsep=0.6cm,footskip=1.2cm]{geometry}
\usepackage[T1]{fontenc}
\usepackage[utf8]{inputenc}
\usepackage{lmodern}
\usepackage{textcomp}
\usepackage[french]{babel}
\usepackage{xcolor}
\usepackage{listings}
\usepackage{graphicx}
\usepackage{array}
\usepackage{enumitem}
\usepackage{amsmath}
\usepackage{tikz}
\usetikzlibrary{positioning,shapes.multipart,arrows.meta}
\usepackage[most]{tcolorbox}

\definecolor{codegreen}{rgb}{0,0.6,0}
\definecolor{codegray}{rgb}{0.5,0.5,0.5}
\definecolor{codepurple}{rgb}{0.58,0,0.82}
\definecolor{backcolour}{rgb}{0.95,0.95,0.92}

\lstdefinestyle{mystyle}{
    backgroundcolor=\color{backcolour},
    commentstyle=\color{codegreen},
    keywordstyle=\color{magenta},
    numberstyle=\tiny\color{codegray},
    stringstyle=\color{codepurple},
    basicstyle=\ttfamily\footnotesize,
    breakatwhitespace=false,
    breaklines=true,
    postbreak=\mbox{\textcolor{codegray}{$\hookrightarrow$}\space},
    captionpos=b,
    keepspaces=true,
    columns=fullflexible,
    numbers=left,
    numbersep=5pt,
    showspaces=false,
    showstringspaces=false,
    showtabs=false,
    upquote=true,
    tabsize=4,
    xleftmargin=1.5em,
    framexleftmargin=0.3em,
    aboveskip=0.6em,
    belowskip=0.6em,
    extendedchars=true,
    literate={é}{{\'e}}1 {è}{{\`e}}1 {ê}{{\^e}}1 {ë}{{\"e}}1 {à}{{\`a}}1 {â}{{\^a}}1
             {ç}{{\c c}}1 {ù}{{\`u}}1 {û}{{\^u}}1 {ô}{{\^o}}1 {î}{{\^i}}1 {ï}{{\"i}}1
             {É}{{\'E}}1 {È}{{\`E}}1 {À}{{\`A}}1 {Ç}{{\c C}}1 {œ}{{\oe}}1
             {«}{{\guillemotleft}}1 {»}{{\guillemotright}}1 {→}{{$\rightarrow$}}1
}
\lstset{style=mystyle, language=Python}
\lstdefinestyle{sqlstyle}{style=mystyle, language=SQL, morekeywords={LIMIT,REFERENCES,IN,TEXT,FLOAT,DATE,VARCHAR,CHAR,REPLACE}, numbers=none, xleftmargin=0.3em, frame=lines}

% Code en ligne
\newcommand{\code}[1]{\texttt{#1}}

% Titre d'une question : numéro et intitulé, barème à droite
\newcommand{\Qt}[2]{\subsection*{#1\hfill{\normalfont\small #2}}}
% Rappel de la consigne en une ligne
\newcommand{\rappel}[1]{\noindent\textit{Ce qui est demandé :} #1\par\smallskip}

% La réponse attendue, mise en évidence
\newtcolorbox{reponse}[1][Réponse]{enhanced, breakable, colback=black!5, colframe=black, boxrule=0.9pt, arc=1mm,
  title={#1}, fonttitle=\bfseries\small, coltitle=black, colbacktitle=black!15, left=2mm, right=2mm, top=1mm, bottom=1mm,
  before skip=0.5em, after skip=0.7em}
% Erreur fréquente
\newtcolorbox{erreur}{enhanced, breakable, colback=white, colframe=black, boxrule=0.5pt, leftrule=3pt, arc=0mm,
  title={Erreur fréquente}, fonttitle=\bfseries\small, coltitle=black, colbacktitle=white, left=2mm, right=2mm, top=1mm, bottom=1mm,
  before skip=0.5em, after skip=0.7em}
% Encadré « À retenir »
\newtcolorbox{retenir}{enhanced, breakable, colback=white, colframe=black, boxrule=0.5pt, arc=1mm, boxsep=1mm,
  borderline={0.4pt}{1.2pt}{black, dashed},
  title={À retenir}, fonttitle=\bfseries\small, coltitle=black, colbacktitle=white, left=2mm, right=2mm, top=1mm, bottom=1mm,
  before skip=0.5em, after skip=0.7em}

\newcommand{\methode}{\noindent\textbf{Méthode.}\ }
\newcommand{\explication}{\noindent\textbf{Explication.}\ }

\setlist{itemsep=0.15em, topsep=0.3em}
\renewcommand{\arraystretch}{1.25}
\setlength{\parindent}{0pt}
\setlength{\parskip}{0.35em}

% Piles dessinées verticalement (sommet en haut) ; \suppr barre un élément supprimé
\newcommand{\suppr}[1]{\ooalign{#1\cr\hidewidth\rule[0.55ex]{1em}{1pt}\hidewidth\cr}}
\newcommand{\pilev}[2]{\begin{tikzpicture}[baseline=(titre.base)]
  \node[font=\scriptsize, align=center, text width=22mm] (titre) at (0,0) {#1};
  \foreach \v [count=\k] in {#2} {\node[draw, minimum width=8mm, minimum height=5mm, inner sep=0pt, font=\small] at (0,-0.25-\k*0.5) {\v};}
\end{tikzpicture}}
\setlength{\emergencystretch}{3em}

\title{Corrigé — CC1 Term NSI - Pile, File et rappels de 1er}
\author{Babash - NSI Tani Malandi et Johanne Bakalara}
\date{Sujet de septembre 2024 (1 h 30, sur 20 points)}

\begin{document}
\pagestyle{headandfoot}
\firstpageheader{}{}{}
\runningheader{\small Corrigé}{}{\small Pile, file et rappels de 1re -- septembre 2024}
\runningheadrule
\firstpagefooter{}{\footnotesize Babash - NSI Tani Malandi et Johanne Bakalara \textperiodcentered{} CC BY-SA 4.0 \textperiodcentered{} page \thepage}{}
\runningfooter{}{\footnotesize Babash - NSI Tani Malandi et Johanne Bakalara \textperiodcentered{} CC BY-SA 4.0 \textperiodcentered{} page \thepage}{}

\maketitle
\thispagestyle{headandfoot}
\begin{center}\footnotesize\itshape Mélange d'un jeu de cartes en deux piles d'après l'exercice 1 du sujet 21-NSIJ2G11 (baccalauréat 2021).\end{center}

\begin{tcolorbox}[colback=white, colframe=black, boxrule=0.5pt, arc=1mm, left=2mm, right=2mm]
\small\textbf{Comment lire ce corrigé.} Pour chaque question : ce qui est demandé, puis la réponse attendue (encadrée), puis l'explication pas à pas. Essaie d'abord seul·e, puis compare ta démarche, pas seulement ton résultat. Tous les codes de ce corrigé ont été exécutés en Python 3.
Le barème est celui du sujet (9 questions, 20 points).
\end{tcolorbox}

\section*{Partie 1 -- Cours}

\Qt{Question 1 -- Questions de cours}{4 points}

\textbf{a)} \rappel{la différence principale entre un tuple et une liste.}
\begin{reponse}
Une \textbf{liste} est \textbf{mutable} : on peut modifier ses éléments, en ajouter ou en retirer après sa création. Un \textbf{tuple} est \textbf{non mutable} (immuable) : une fois créé, on ne peut plus le modifier.
\end{reponse}
\explication On les reconnaît aussi à l'écriture : une liste s'écrit entre crochets, \code{[1, 2, 3]}, un tuple entre parenthèses, \code{(1, 2, 3)}. Exemple :
\begin{lstlisting}
liste = [1, 2, 3]
liste[0] = 10        # possible : liste vaut [10, 2, 3]
couple = (1, 2, 3)
couple[0] = 10       # TypeError: 'tuple' object does not support item assignment
\end{lstlisting}
Le message d'erreur signifie « un objet de type tuple n'accepte pas l'affectation d'un élément ».

\textbf{b)} \rappel{l'instruction qui donne le numéro de Bob dans \code{repertoire}.}
\begin{reponse}
\code{repertoire[1]["Numero"]} (cette expression vaut \code{"0600000002"}).
\end{reponse}
\explication \code{repertoire} est une \textbf{liste de dictionnaires}. On lit l'expression de gauche à droite :
\begin{enumerate}
\item \code{repertoire[1]} : l'élément d'indice 1 de la liste (le deuxième, car on compte à partir de 0), c'est le dictionnaire de Bob, \code{\{"Prenom": "Bob", "Numero": "0600000002"\}} ;
\item \code{["Numero"]} : dans ce dictionnaire, la valeur associée à la clé \code{"Numero"}.
\end{enumerate}

\begin{erreur}
Écrire \code{repertoire["Bob"]} : une liste s'indexe par un nombre, pas par un texte. Python affiche \code{TypeError: list indices must be integers or slices, not str} (« les indices de liste doivent être des entiers »).
\end{erreur}

\textbf{c)} \rappel{la structure de données qui correspond à LIFO.}
\begin{reponse}
La \textbf{pile}. LIFO veut dire \textit{Last In, First Out} : « dernier entré, premier sorti ».
\end{reponse}
\explication Dans une pile, on ajoute et on retire au même endroit, le sommet : le dernier élément ajouté est le premier retiré. La file est FIFO (\textit{First In, First Out}, « premier entré, premier sorti »).

\textbf{d)} \rappel{les boîtes rangées les unes derrière les autres : pile ou file ?}
\begin{reponse}
C'est une \textbf{pile} : la première boîte mise en rayon est la dernière prise.
\end{reponse}
\explication La première boîte posée se retrouve tout au fond ; les clientes et les clients prennent la boîte de devant, qui est la dernière posée. « Premier entré, dernier sorti », c'est la même chose que « dernier entré, premier sorti » : c'est une pile. Pour une file, il faudrait remplir le rayon par l'arrière et servir par l'avant.

\section*{Partie 2 -- Compréhension de code}

\begin{lstlisting}
def mystere(tab):
    i = 0
    j = len(tab) - 1
    while i != j :
        if tab[i] == 0:
            i = i + 1
        else :
            valeur = tab[j]
            tab[j] = 1
            tab[i] = valeur
            j= j - 1
    return tab
\end{lstlisting}

\Qt{Question 2 -- Trace de \code{mystere([0, 1, 0, 1, 0, 1])}}{2 points}
\rappel{donner \code{i}, \code{j} et \code{tab} à chaque tour de boucle.}

\methode
\begin{enumerate}
\item Noter les valeurs avant la boucle : \code{i = 0}, \code{j = len(tab) - 1 = 5}.
\item À chaque tour, regarder \code{tab[i]} : si c'est 0, on avance \code{i} ; sinon, on met la valeur de \code{tab[j]} dans la case \code{i}, un 1 dans la case \code{j}, et on recule \code{j}.
\item S'arrêter quand \code{i} et \code{j} sont égaux.
\end{enumerate}

\begin{reponse}
\begin{tabular}{|c|c|c|c|l|}\hline
Tour & \code{tab[i]} lu & \code{i} & \code{j} & \code{tab} en fin de tour \\\hline
avant la boucle & & 0 & 5 & \code{[0, 1, 0, 1, 0, 1]} \\\hline
1 & \code{tab[0]} = 0 & 1 & 5 & \code{[0, 1, 0, 1, 0, 1]} \\
2 & \code{tab[1]} = 1 & 1 & 4 & \code{[0, 1, 0, 1, 0, 1]} \\
3 & \code{tab[1]} = 1 & 1 & 3 & \code{[0, 0, 0, 1, 1, 1]} \\
4 & \code{tab[1]} = 0 & 2 & 3 & \code{[0, 0, 0, 1, 1, 1]} \\
5 & \code{tab[2]} = 0 & 3 & 3 & \code{[0, 0, 0, 1, 1, 1]} \\\hline
\end{tabular}\\
Après le tour 5, \code{i} et \code{j} valent 3 : la boucle s'arrête et la fonction renvoie \code{[0, 0, 0, 1, 1, 1]}.
\end{reponse}

\explication
\begin{itemize}
\item Tour 2 : \code{tab[1]} vaut 1. On copie \code{tab[5]} (qui vaut 1) dans \code{valeur}, on met 1 dans \code{tab[5]} et \code{valeur} dans \code{tab[1]}. Les deux cases valaient déjà 1 : le tableau ne change pas, mais \code{j} recule à 4.
\item Tour 3 : \code{tab[1]} vaut toujours 1. Cette fois, \code{tab[4]} vaut 0 : le 0 passe en case 1 et le 1 en case 4. On obtient \code{[0, 0, 0, 1, 1, 1]} et \code{j} recule à 3.
\item Remarque : \code{i} n'avance pas quand \code{tab[i]} vaut 1 ; on revérifie la même case au tour suivant.
\end{itemize}

\begin{erreur}
Avancer \code{i} à chaque tour, comme dans une boucle \code{for}. Ici, \code{i} n'avance que si \code{tab[i]} vaut 0, et \code{j} ne recule que dans l'autre cas : à chaque tour, une seule des deux variables change.
\end{erreur}

\Qt{Question 3 -- Réfléchir sur la fonction}{2 points}

\textbf{a)} \rappel{la valeur renvoyée pour quatre appels.}
\begin{reponse}
\begin{tabular}{@{}ll@{}}
\code{mystere([0, 0, 0])} & renvoie \code{[0, 0, 0]} \\
\code{mystere([])} & provoque une erreur \code{IndexError} \\
\code{mystere([0, 0, 1])} & renvoie \code{[0, 0, 1]} \\
\code{mystere([1, 1, 1, 1])} & renvoie \code{[1, 1, 1, 1]} \\
\end{tabular}
\end{reponse}
\explication
\begin{itemize}
\item \code{[0, 0, 0]} : \code{j} vaut 2 ; \code{i} avance 0, 1, 2 et la boucle s'arrête. Rien n'a changé.
\item \code{[]} : \code{j} vaut \code{len([]) - 1}, soit $-1$. Comme \code{i = 0} est différent de $-1$, on entre dans la boucle et on lit \code{tab[0]}, qui n'existe pas : erreur « indice de liste hors limites ». La boucle ne renvoie donc rien pour un tableau vide.
\item \code{[0, 0, 1]} : \code{i} avance jusqu'à 2, où il rejoint \code{j} : le tableau est inchangé.
\item \code{[1, 1, 1, 1]} : \code{tab[0]} vaut toujours 1, donc \code{j} recule 3, 2, 1, 0 et rejoint \code{i} : le tableau est inchangé.
\end{itemize}

\textbf{b)} \rappel{un nom plus parlant pour la fonction.}
\begin{reponse}
Par exemple \code{trier\_zeros\_uns} (ou \code{ranger\_0\_avant\_1}) : la fonction trie un tableau de 0 et de 1, en plaçant tous les 0 au début et tous les 1 à la fin.
\end{reponse}
\explication Pour trouver ce que fait une fonction, compare les entrées et les sorties de plusieurs essais : \code{[0, 1, 0, 1, 0, 1]} devient \code{[0, 0, 0, 1, 1, 1]}, et un tableau déjà rangé ne change pas. \code{i} parcourt les 0 du début, \code{j} remplit la fin avec des 1.

\section*{Partie 3 -- Rédaction d'un code simple}

\Qt{Question 4 -- La fonction \code{moyenne}}{4 points}
\rappel{renvoyer la moyenne des valeurs d'un tableau non vide de flottants.}

\methode
\begin{enumerate}
\item Créer une variable \code{somme} qui vaut 0 au départ.
\item Parcourir le tableau et ajouter chaque valeur à \code{somme}.
\item Diviser la somme par le nombre de valeurs, \code{len(tab)}, et renvoyer le résultat.
\end{enumerate}

\begin{reponse}
\begin{lstlisting}
def moyenne(tab):
    somme = 0
    for valeur in tab:
        somme = somme + valeur
    return somme / len(tab)
\end{lstlisting}
\end{reponse}

\explication
\begin{itemize}
\item Ligne 2 : on initialise l'\textbf{accumulateur} \code{somme} (une variable qui « accumule » les valeurs).
\item Lignes 3 et 4 : la boucle prend chaque valeur du tableau, l'une après l'autre, et l'ajoute à \code{somme}.
\item Ligne 5 : le \code{return} est \textbf{après} la boucle (il n'est pas indenté sous le \code{for}). Le tableau n'est pas vide, donc on ne divise jamais par 0.
\end{itemize}

Trace pour \code{moyenne([1.0, 2.0, 4.0])} :
\begin{center}\small
\begin{tabular}{|c|c|c|}\hline
Tour & \code{valeur} & \code{somme} en fin de tour \\\hline
avant & & 0 \\
1 & 1.0 & 1.0 \\
2 & 2.0 & 3.0 \\
3 & 4.0 & 7.0 \\\hline
\end{tabular}
\end{center}
La fonction renvoie $7{,}0 / 3$, que Python affiche \code{2.3333333333333335} (un flottant n'a qu'un nombre limité de chiffres, d'où le 5 à la fin). Et \code{moyenne([1.0])} renvoie \code{1.0}.

\textbf{Variante} avec la fonction prédéfinie \code{sum} : \code{return sum(tab) / len(tab)}. Elle est juste, mais au contrôle on attend souvent la boucle, qui montre que tu sais calculer une somme.

\begin{erreur}
Placer le \code{return} dans la boucle (indenté sous le \code{for}) : la fonction s'arrête dès le premier tour et renvoie \code{1.0 / 3} au lieu de la moyenne. Autre erreur : oublier de diviser, ou diviser par un nombre écrit à la main (\code{/ 3}) au lieu de \code{len(tab)}.
\end{erreur}

\section*{Partie 4 -- Exercice type bac : piles}

\begin{retenir}
Une \textbf{pile} n'a ici que quatre opérations : \code{Pile()} crée une pile vide, \code{p.est\_vide()} teste si elle est vide, \code{p.empiler(e)} pose \code{e} au sommet, \code{p.depiler()} retire l'élément du sommet et le renvoie. On ne peut pas lire un élément au milieu de la pile.
\end{retenir}

\Qt{Question 5 -- Mélanger \code{['7','8','9','10','V','D','R','A']}}{1 point}
\rappel{la liste obtenue après le mélange.}
\begin{reponse}
\code{['10', 'A', '9', 'R', '8', 'D', '7', 'V']}
\end{reponse}
\explication On suit les étapes de l'énoncé, avec $N = 8$, donc $N/2 = 4$ :
\begin{enumerate}
\item Pile de gauche : les 4 premiers éléments \code{'7'}, \code{'8'}, \code{'9'}, \code{'10'}, empilés dans cet ordre ; le sommet est \code{'10'}.
\item Pile de droite : \code{'V'}, \code{'D'}, \code{'R'}, \code{'A'} ; le sommet est \code{'A'}.
\item On prend en alternance le sommet de gauche, puis celui de droite : \code{'10'}, \code{'A'}, puis \code{'9'}, \code{'R'}, puis \code{'8'}, \code{'D'}, puis \code{'7'}, \code{'V'}.
\end{enumerate}
\begin{center}
\pilev{pile de gauche}{'10','9','8','7'}\qquad\qquad
\pilev{pile de droite}{'A','R','D','V'}
\end{center}
On vérifie la méthode sur l'exemple de l'énoncé : \code{['V','D','R','3','7','10']} donne bien \code{['R','10','D','7','V','3']}.

\Qt{Question 6 -- La fonction \code{liste\_vers\_pile}}{2 points}
\rappel{compléter la fonction pour que le dernier élément de la liste soit au sommet de la pile.}
\begin{reponse}
\begin{lstlisting}
def liste_vers_pile(L):
    '''prend en paramètre une liste et renvoie une pile'''
    N = len(L)
    p_temp = Pile()
    for i in range(N):
        p_temp.empiler(L[i])
    return p_temp
\end{lstlisting}
\end{reponse}
\explication On parcourt la liste du premier au dernier élément (\code{i} va de 0 à $N-1$, car \code{range(N)} s'arrête avant $N$) et on empile chaque élément. Le dernier empilé, \code{L[N-1]}, est donc au sommet, comme demandé. On renvoie la pile, pas la liste.

Exemple : \code{liste\_vers\_pile([1, 2, 3])} renvoie une pile dont le sommet est 3 ; un premier \code{depiler()} renvoie 3.

\begin{erreur}
Écrire \code{p\_temp.empiler(i)} : on empile alors les \textbf{indices} 0, 1, 2… au lieu des \textbf{éléments} \code{L[0]}, \code{L[1]}, \code{L[2]}…
\end{erreur}

\Qt{Question 7 -- Affichages de \code{partage([1,2,3,4,5,6])}}{2 points}
\rappel{ce qu'affichent les deux appels à \code{affichage\_pile}.}
\begin{reponse}
\begin{minipage}[t]{0.45\linewidth}
Premier affichage (\code{p\_gauche}) :\\
\code{3}\\\code{2}\\\code{1}
\end{minipage}\hfill
\begin{minipage}[t]{0.45\linewidth}
Second affichage (\code{p\_droite}) :\\
\code{6}\\\code{5}\\\code{4}
\end{minipage}\\[0.4em]
(La pile est affichée verticalement, sommet en haut.)
\end{reponse}
\explication
\begin{enumerate}
\item \code{N} vaut 6, donc \code{N // 2} vaut 3 (\code{//} est la division entière).
\item Première boucle : \code{i} prend les valeurs 0, 1, 2 ; on empile \code{L[0]}, \code{L[1]}, \code{L[2]}, c'est-à-dire 1, 2, 3 dans \code{p\_gauche}. Le sommet est 3.
\item Seconde boucle : \code{range(3, 6)} donne 3, 4, 5 ; on empile 4, 5, 6 dans \code{p\_droite}. Le sommet est 6.
\item Chaque pile s'affiche de haut en bas : d'abord 3, 2, 1, puis 6, 5, 4.
\end{enumerate}

\Qt{Question 8 -- Expliquer la fusion de deux piles}{1 point}
\rappel{expliquer, avec des schémas, comment fusionner \code{p\_gauche} et \code{p\_droite} en une liste en alternant.}
\begin{reponse}
On part d'une liste vide \code{L}. Tant que les piles ne sont pas vides : on dépile le sommet de \code{p\_gauche} et on l'ajoute à la fin de \code{L}, puis on dépile le sommet de \code{p\_droite} et on l'ajoute à la fin de \code{L}. Quand les piles sont vides, \code{L} contient tous les éléments, en alternant gauche et droite.
\end{reponse}
\explication Les schémas, avec les piles de la question 7 (sommet en haut) :
\begin{center}\small
\begin{tabular}{|c|c|c|l|}\hline
Étape & \code{p\_gauche} & \code{p\_droite} & \code{L} \\\hline
départ & 3 2 1 & 6 5 4 & \code{[]} \\
1 & 2 1 & 5 4 & \code{[3, 6]} \\
2 & 1 & 4 & \code{[3, 6, 2, 5]} \\
3 & vide & vide & \code{[3, 6, 2, 5, 1, 4]} \\\hline
\end{tabular}
\end{center}
Les deux piles ont le même nombre d'éléments ($N/2$ chacune), donc elles se vident en même temps.

\Qt{Question 9 -- La fonction \code{fusion}}{2 points}
\rappel{écrire \code{fusion(p1, p2)} qui renvoie la liste construite à partir des deux piles.}
\begin{reponse}
\begin{lstlisting}
def fusion(p1, p2):
    L = []
    while not p1.est_vide():
        L.append(p1.depiler())
        L.append(p2.depiler())
    return L
\end{lstlisting}
\end{reponse}
\explication
\begin{itemize}
\item Ligne 2 : la liste résultat est vide au départ.
\item Ligne 3 : on répète tant que \code{p1} n'est pas vide (les deux piles ont la même taille, il suffit d'en tester une). On utilise \code{while} et non \code{for}, car on ne peut pas connaître la taille d'une pile avec les seules opérations de l'énoncé.
\item Lignes 4 et 5 : \code{depiler()} retire le sommet et le renvoie ; \code{append} l'ajoute à la fin de la liste. On prend d'abord dans \code{p1}, puis dans \code{p2}.
\end{itemize}
Exemple : avec les piles de \code{partage(['V','D','R','3','7','10'])}, \code{fusion} renvoie \code{['R', '10', 'D', '7', 'V', '3']}, le résultat de l'énoncé.

\begin{erreur}
Oublier les parenthèses : \code{p1.est\_vide} (sans \code{()}) n'appelle pas la méthode. C'est un objet « méthode », toujours considéré comme vrai : \code{not p1.est\_vide} vaut \code{False} et la boucle ne s'exécute jamais. La fonction renvoie alors une liste vide.
\end{erreur}

\end{document}
